% ============================================================================
% EVIDENCE BEFORE ERASURE -- Research report scaffold (Track 2)
% ============================================================================
%
% READ THIS BLOCK BEFORE YOU SUBMIT ANYTHING.
%
% 1. THIS IS SCAFFOLDING, NOT A FINISHED REPORT. Every bullet list in this file
%    is fact-dense NOTES pulled from the project's own audit trail (with a
%    citation to where each fact lives), not finished prose. The sprint's own
%    rules say a report "may be penalized if it reads as generated rather than
%    written by the team" (padded sections, generic framing, claims without
%    sources). Turn these bullets into your own sentences. Do not submit this
%    file's bullet lists as flowing paragraphs verbatim.
%
% 2. THE OFFICIAL TEMPLATE IS A GOOGLE DOC, not LaTeX:
%    https://docs.google.com/document/d/1PQBlhI3tM5vb51x7jBWXBQMYg6hkiU_x8RaCws4kjl4/copy?usp=sharing
%    The FAQ says to always use whatever is linked on the Guidelines tab (the
%    copy in acceptance emails may be stale). This .tex file matches the
%    REQUIRED STRUCTURE AND LIMITS (8 pages excl. refs/appendices, abstract
%    <=150 words, required Limitations & Dual-Use appendix, author/affiliation
%    block) -- but before final submission, either port this content into the
%    actual Google Doc template, or confirm a LaTeX PDF is acceptable in this
%    round's Guidelines tab.
%
% 3. RESULTS ARE NOT YET REAL. X13 (the actual experiment) has not been run as
%    of this draft -- only Gate 1 (engineering) and Gate 2 (independent
%    readiness) have passed. Every number in \S Results is marked PENDING.
%    Do not invent numbers to fill the gap; an honest "results pending, here is
%    exactly what will be reported and how" is a legitimate, common thing to
%    submit under time pressure, and is far better than a fabricated number.
%
% 4. FILL IN: author names/affiliations, the exact title (a working title is
%    given below), and every [TODO] marker.
%
% Compile with pdflatex (no exotic packages; should work on Overleaf or a
% local TeXLive install unmodified).
% ============================================================================

\documentclass[11pt]{article}

\usepackage[margin=1in]{geometry}
\usepackage{times}
\usepackage{booktabs}
\usepackage{enumitem}
\usepackage{hyperref}
\usepackage{graphicx}
\usepackage{caption}
\usepackage[table]{xcolor}
\usepackage{titlesec}
\usepackage{fancyhdr}

\hypersetup{
  colorlinks=true,
  linkcolor=blue!50!black,
  citecolor=blue!50!black,
  urlcolor=blue!50!black
}

\titlespacing*{\section}{0pt}{12pt}{6pt}
\titlespacing*{\subsection}{0pt}{10pt}{4pt}

\pagestyle{fancy}
\fancyhf{}
\fancyhead[L]{\small Evidence Before Erasure --- Draft, not for distribution}
\fancyhead[R]{\small \thepage}
\renewcommand{\headrulewidth}{0.2pt}

% Simple TODO/pending markers, visible in the compiled PDF so nobody forgets
% to resolve them before submission.
\newcommand{\todo}[1]{\textcolor{red}{\textbf{[TODO: #1]}}}
\newcommand{\pending}[1]{\textcolor{orange!80!black}{\textbf{[PENDING -- #1]}}}

\title{\vspace{-1.5em}
  Evidence Before Erasure: Does Event-Driven Monitoring Preserve\\
  More Forensic Evidence of AI Agent Actions Than Periodic Polling?\\
  \todo{confirm/replace working title}}

\author{
  \todo{Author 1 Name}$^{1}$ \quad
  \todo{Author 2 Name}$^{1}$ \quad
  \todo{Author 3 Name}$^{1}$ \quad
  \todo{Author 4 Name}$^{1}$ \\[4pt]
  \small $^{1}$\todo{affiliation(s)} \\
  \small \texttt{\todo{contact email(s)}}
}
\date{\today}

\begin{document}
\maketitle
\thispagestyle{fancy}

\begin{abstract}
\noindent
% ABSTRACT -- HARD LIMIT 150 WORDS. Count before submitting.
% Draft below is ~140 words; trim/rewrite in your own words.
When an AI agent takes an action inside a shared, mutable environment
(edits a page, coordinates with other agents, deletes evidence of its own
behavior), the trace of that action has a finite shelf life: it can be
overwritten, deleted, or evicted from a storage-capped monitor before anyone
looks. We study, in a trace-grounded simulation of a wiki-like environment
populated by autonomous agents, whether a simple time-interval polling
monitor retains less \emph{critical forensic evidence} of agent actions than
a budget-matched, event-driven monitor under identical storage and request
constraints. We built and independently, adversarially audited a 65-item
human-labeled benchmark of concrete agent behaviors (including deliberate
evasion patterns such as backup-page creation before deletion) and a
retained-only evaluator with a formally specified causal-ordering model.
\pending{one-sentence headline result once X13 has run}.
\end{abstract}

\vspace{0.5em}
\noindent\textbf{Keywords:} \todo{3--5 keywords, e.g. AI incident response,
observability, evidence retention, monitoring policy, simulation}

% ==========================================================================
\section{Introduction}
% ==========================================================================
% Required content per the sprint rubric: which track/sub-problem, why it
% matters, what the artifact is for. Write 3-5 short paragraphs from these
% points -- do not just list them.

\begin{itemize}[leftmargin=*]
  \item \textbf{Track and problem.} Track 2 (AI Incident Response). The
    sub-problem: incident response depends on evidence surviving long enough
    to be collected; monitoring/logging design choices (poll interval,
    storage cap, feed- vs. schedule-driven collection) determine what
    evidence is even \emph{retrievable} after the fact, independent of
    whether the incident was ever noticed in time.
  \item \textbf{Why it matters.} This reframes ``AI incident response'' as
    partly an information-preservation/observability engineering problem, not
    only a detection problem. If a monitoring architecture systematically
    loses certain classes of evidence (e.g., self-erasing or fast-mutating
    agent behavior), post-hoc investigation of an incident can fail even when
    the incident is eventually flagged --- the evidence needed to explain it
    is already gone. \todo{cite 1-2 real incident-response / observability
    sources establishing that evidence retention -- not just detection -- is
    a recognized bottleneck; the sprint rubric penalizes claims without
    sources}.
  \item \textbf{Concrete motivating example (use this -- it is real, already
    independently re-verified, and concrete).} In the simulated corpus, an
    agent (handle \texttt{Aug17ConstructionAgent}) posted a warning that a
    "wiki cleanup/deletion sweep appears active," naming a specific backup
    page to use if the original vanished. 96 seconds later that backup page
    (\texttt{ZZZDataUSAConstructionWageLive}) was created. The original page
    and the backup page were then deleted 12 seconds apart. Every timestamp
    and hash in this example was independently re-derived from the raw
    revision/event export, not copied from a prior summary
    (\texttt{docs/JUDGE\_WALKTHROUGH.md}). This is exactly the shape of
    evidence a badly-timed or too-infrequent monitor would miss entirely.
  \item \textbf{What the artifact is.} (a) A frozen, hand-adjudicated
    benchmark of 65 ``critical'' agent-behavior propositions grounded in a
    13,000+-revision simulated wiki trace, each requiring exact retained
    textual evidence (and, for 8 of the 65, additional context); (b) a
    retained-only evaluator with a formally specified causal-ordering model
    (an operation clock that provably agrees between a full run and a
    truncated ``fresh prefix'' up to any checkpoint); (c) a comparison
    harness for periodic (P/PD/PCD at several intervals) vs.\ event-driven
    (token-bucket-rate-limited) collection, plus four adversarial ``hostile
    controls'' (a low-information baseline, an online repair variant, a
    reversed-service-order stress test, and a terminal archive-augmentation
    control) built specifically to stress-test whether any observed
    difference is an artifact of an unfair comparison.
  \item \textbf{Scope discipline (worth stating explicitly in the report --
    it is unusual and a strength, not a hedge).} The primary contrast, the
    benchmark, the denominator (held-out critical core, $K=23$), and the
    decision threshold were frozen \emph{before} any real result was
    computed; every methodological repair found during audit was applied
    outcome-blind. \todo{state plainly, once results exist, whether any
    parameter was changed after seeing a result --- if none were, say so
    explicitly, it is a real strength}.
\end{itemize}

% ==========================================================================
\section{Related Work}
% ==========================================================================
% Required: what prior work this builds on, what it already does, what gap
% remains, what is new here. \todo -- this section needs real citations the
% team gathers; do not let an AI tool invent citations.

\begin{itemize}[leftmargin=*]
  \item \todo{Cite prior AI incident/agent-monitoring work this builds on
    (e.g. agent evaluation/red-teaming benchmarks, logging/observability
    literature, digital forensics evidence-retention literature, prior work
    on periodic vs.\ event-driven monitoring in distributed systems more
    generally).}
  \item \todo{State plainly what existing benchmarks/evaluations already
    measure (typically: whether an agent's actions are \emph{detected} or
    \emph{scored} for harmfulness) and what they do \emph{not} measure
    (whether the evidence needed to investigate an action \emph{survives}
    under realistic storage/collection constraints).}
  \item \textbf{Gap and contribution (draft, tighten with citations):} to our
    knowledge, no existing AI incident-response evaluation treats evidence
    \emph{retention under a capacity-constrained monitor} as the primary
    outcome variable, or builds an adversarial ``hostile control'' suite
    specifically to stress-test the fairness of a monitoring-policy
    comparison. \todo{soften/qualify this claim once related work is actually
    checked -- an unverified novelty claim is exactly what loses
    Impact/Innovation points with judges.}
\end{itemize}

% ==========================================================================
\section{Methodology}
% ==========================================================================
% Required: enough detail to replicate -- methods, experimental design, data,
% sources, assumptions, evaluation procedure. This project's own docs already
% contain nearly all of this in verified form; the bullets below cite exactly
% where. Turn into prose; keep the specific numbers, drop anything not needed
% for replication understanding at report length.

\subsection{Environment and data}
\begin{itemize}[leftmargin=*]
  \item Trace-grounded simulation over a released wiki-style corpus (``DSE'')
    of autonomous-agent edits: \todo{state corpus size --- 13{,}403+ revisions
    per \texttt{docs/PRE\_RESULTS\_CLOSURE.md}; confirm exact final count} across
    a fixed replay window (\texttt{2026-05-24} to \texttt{2026-07-15}).
  \item Study description (use verbatim or close paraphrase --- this exact
    framing was deliberately frozen as the paper's own epistemic-status
    statement, per \texttt{docs/X13\_RUN\_CONTRACT.md}): \emph{``Trace-grounded,
    live-only simulation under a shared hypothetical change-feed interface,
    conditional on the released DSE history. Evidence not retained by a
    simulated collector is not evidence permanently erased from the
    website.''}
\end{itemize}

\subsection{Benchmark}
\begin{itemize}[leftmargin=*]
  \item 65 hand-labeled propositions (concrete claims that an agent did
    something), 57 critical / 8 non-critical; split 38 development / 27
    held-out, 34/23 critical respectively. Primary held-out critical
    denominator $K = 23$.
  \item 25 title-groups, 52 catalog pages including context dependencies;
    1{,}421 core occurrence (span-provenance) rows across the corpus.
  \item 8 of 65 propositions require additional \emph{context} (a prior
    revision or a deletion event) beyond their own anchor span to be
    meaningful; all 8 have frozen, groundable context alternatives (57 are
    self-contained). One proposition
    (\texttt{PROP-20260616-61}) required adjudicating a genuine three-way
    ambiguity (three near-identical candidate ``older helper pages''
    matching a referenced query pattern) rather than forcing a single
    choice --- resolved as an OR across all three valid alternatives.
  \item Every quotation, span, and hash in the benchmark is independently
    re-derivable from the raw export by construction (SHA-256 over the
    reconstructed source bytes, exact half-open character spans).
  \item \textbf{Novelty/replication note worth including:} the annotation
    process itself was adversarially audited across the team multiple times;
    concrete, disclosed corrections found along the way included a stale
    benchmark count in an early spec, an encoding bug that silently mismatched
    2 evidence rows and 80 occurrence rows against their true source bytes
    (traced to an unconditional Latin-1 byte-projection requirement in the
    raw export, not a declared-encoding branch), and a first-pass context
    axis that could initially count context without core. \todo{decide how
    much of this process detail belongs in the main report vs.\ an appendix
    --- it supports Execution Quality but can bloat the 8-page limit.}
\end{itemize}

\subsection{Collection policies compared}
\begin{itemize}[leftmargin=*]
  \item \textbf{Periodic:} fixed-interval polling at several intervals
    $\Delta$ (P, PD, PCD variants), evaluated across all 60 phase offsets
    within each interval (not a single arbitrarily-chosen phase).
  \item \textbf{Event-driven (E):} feed-triggered collection, rate-limited by
    a token-bucket parameter $q$.
  \item \textbf{Primary contrast (frozen before any result was seen):} PCD at
    $\Delta=15$ minutes (mean across all 60 phases) vs.\ E at $q=30$; shared
    1\,MiB capped, deduplicated store; 5s feed lag; 60s poll interval; 0s GET
    delay. Decision rule: a $\geq 10$ percentage-point difference in held-out
    critical core retention at the study endpoint.
  \item \textbf{Hostile controls,} run to test whether the primary contrast
    is a fair/robust comparison rather than an artifact of policy choice:
  \begin{itemize}
    \item \emph{F} --- a low-information baseline: one terminal-state
      directory listing plus one atomic GET per currently-live title, no
      ongoing polling at all.
    \item \emph{PCD-R} --- an online self-repair variant of periodic
      collection that opportunistically re-fetches known-stale content,
      including under non-zero response delay with a real pending-request
      model (a request already in flight for a title suppresses a duplicate
      repair attempt for that same title).
    \item \emph{Reverse-order} stress --- services titles/events in reversed
      order at each sweep, to test whether the comparison is sensitive to
      arbitrary service ordering under storage pressure.
    \item \emph{Terminal archive augmentation} --- a full enumeration of every
      revision saved by the study endpoint, layered on top of the live
      collection, to bound how much a terminal archival pass alone recovers.
  \end{itemize}
\end{itemize}

\subsection{Evaluator and causal-ordering model}
\begin{itemize}[leftmargin=*]
  \item The evaluator scores \emph{only} what a policy's own capped,
    deduplicated store actually retained at the study checkpoint --- never a
    collector's full diagnostic/rejected-body history. A proposition counts
    as covered only if a complete, exactly-hashed, exactly-spanned frozen
    occurrence is present in the retained snapshot at that checkpoint (OR
    across independently sufficient alternatives, AND within a multi-span
    alternative).
  \item \textbf{Causal-ordering guarantee (worth stating as a concrete,
    testable property, not just an assumption):} a single monotonic
    operation clock governs feed polls, GET dispatch, and completions; an
    as-of-checkpoint cost/coverage query computed from a full run is required
    to exactly equal the same query computed by truncating the run to a
    fresh prefix ending at that checkpoint. This was verified, not assumed
    (independently constructed adversarial histories, not just the
    project's own unit tests).
  \item Unknown/unavailable required context is retained as an explicit
    interval (lower/upper bound), never coerced to a zero or a point
    estimate.
  \item \todo{once real numbers exist: state exact software/commit identity
    used for the reported run (git commit hash, frozen manifest hash) so the
    result is exactly reproducible.}
\end{itemize}

\subsection{Independent verification process}
\begin{itemize}[leftmargin=*]
  \item Every methodological/engineering component passed through an
    explicit two-gate process before any real scoring was authorized: an
    engineering-readiness gate, then a comprehensive independent audit gate
    (adversarial, hand-derived expected values written down \emph{before}
    executing code, on fixtures the auditor constructed independently rather
    than reusing the implementer's own test fixtures).
  \item \todo{state final gate outcome plainly once locked for submission:
    both gates passed as of this draft; no real scoring has occurred yet.}
  \item This process caught and fixed concrete implementation bugs before
    any result existed, e.g.\ a data-integrity check that failed open (raised
    an untyped error rather than a specific, catchable integrity violation)
    on a corrupted retained-body reference, and a real-mode cost/overhead
    reporting gap where several required per-row measurements were never
    populated. \todo{decide exact phrasing/level of detail for the report ---
    this is genuine Execution Quality evidence but should be summarized, not
    itemized at full length.}
\end{itemize}

% ==========================================================================
\section{Results}
% ==========================================================================
\pending{This entire section. X13 (the actual 791-configuration experiment)
has not been executed as of this draft. Do not fill in numbers here until a
real, authorized run has actually completed. Below is the exact shape the
results section will take once it has.}

\subsection{Primary result}
\begin{itemize}[leftmargin=*]
  \item \pending{Held-out critical core retention: PCD15 (mean of 60 phases)
    vs.\ E30, percentage-point difference, and which side of the frozen
    $\geq 10$pp threshold it falls on.}
  \item \pending{Table: per-phase PCD15 retention distribution (all 60
    phases: min/median/mean/p10/p90) vs.\ single E30 value.}
\end{itemize}

\subsection{Sensitivity / hostile-control results}
\begin{itemize}[leftmargin=*]
  \item \pending{F baseline retention (lower bound / sanity floor).}
  \item \pending{PCD-R vs.\ plain PCD --- does online repair close any gap?}
  \item \pending{Reverse-order stress --- does service order change the
    outcome materially?}
  \item \pending{Archive augmentation --- how much does a terminal archival
    pass alone recover, live-only vs.\ archive-augmented?}
  \item \pending{Cost/overhead comparison table --- requests, bytes,
    byte-hours, wall-clock/CPU/RSS --- at matched access, per the frozen
    accounting contract.}
\end{itemize}

\subsection{Main threat to validity}
\begin{itemize}[leftmargin=*]
  \item \todo{fill in honestly once results exist --- candidates already
    identified during methodology design: (a) benchmark is 65
    hand-adjudicated items, not a large statistical sample; (b) simulated
    environment, not live production telemetry; (c) sampling/episode
    provenance for the underlying corpus is not independently certified as
    representative (explicitly flagged during the project's own audit).}
\end{itemize}

% ==========================================================================
\section{Discussion}
% ==========================================================================
\begin{itemize}[leftmargin=*]
  \item \todo{What the result implies for monitoring-policy design in AI
    incident response, once it exists --- write this only after Results is
    real; do not pre-write a conclusion.}
  \item \textbf{Practical implication regardless of which way the result
    goes (safe to write now):} the benchmark, evaluator, and hostile-control
    suite are themselves reusable artifacts for testing any future monitoring
    policy against the same frozen evidence-retention question, on this or
    an extended corpus.
  \item \todo{Future work: e.g.\ broader capacity-cap sweep (explicitly
    scoped out of this run as \texttt{NOT\_EXECUTED\_SCOPE\_AMENDMENT}), real
    (non-simulated) agent telemetry, additional collection policies.}
\end{itemize}

% ==========================================================================
\section*{Acknowledgment of AI Tool Use}
% ==========================================================================
% The sprint explicitly requires disclosure of what was built on / what AI
% assistance was used for, distinct from "did AI write the report." Be
% specific and honest here -- this protects the team, it does not hurt the
% Impact/Innovation score.
\todo{State specifically: which AI tools were used, for what (e.g.\ code
implementation assistance, independent adversarial code review, debugging),
and confirm this report's prose was written by the team, with AI used only
to check reasoning/find gaps, per the sprint's own stated policy on AI tool
use in \texttt{docs/Hackathon\_Judging\_Submission\_Requirements.md}.}

% ==========================================================================
\appendix
\section{Limitations and Dual-Use Considerations}
% ==========================================================================
% REQUIRED APPENDIX. Do not omit. Write in full sentences; bullets below are
% starting points only, drawn from things already explicitly disclosed
% during this project's own audit trail (cite the internal doc each time so
% every claim has a source, per the sprint's evidence-based-claims rule).

\subsection*{Limitations}
\begin{itemize}[leftmargin=*]
  \item \textbf{Simulated environment, not production telemetry.} All agent
    behavior is drawn from a released simulated wiki trace, not real deployed
    AI agents; findings characterize retention dynamics under this
    trace-and-model, not a general claim about production monitoring
    systems.
  \item \textbf{Benchmark scale and provenance.} 65 hand-adjudicated
    propositions is not a large statistical sample; the underlying corpus's
    sampling/episode provenance is explicitly not certified as a
    representative, exhaustively-validated incident population (a
    limitation the team identified and disclosed during its own
    methodology audit, not one an external reviewer found).
  \item \textbf{Scope amendment.} A originally-larger experimental matrix was
    narrowed under time constraints (e.g.\ a broader storage-cap sweep and
    several secondary cross-products were explicitly descoped rather than
    silently dropped); every omission is logged with an explicit reason
    rather than presented as completed work. \todo{state exact omitted items
    if space allows, or reference the scope-amendment record.}
  \item \textbf{No claim of real-world permanent erasure.} A body missing
    from a simulated collector's retained store is not equivalent to
    evidence being permanently destroyed on the real underlying
    system --- only that this particular simulated monitor did not retain
    it.
  \item \textbf{No claim about agent intent.} Wiki actor handles in the
    corpus are labels, not independently verified distinct agent identities
    or confirmed intentional evasion; behavioral patterns that look like
    evasion (e.g.\ backup-page creation before a deletion) are reported as
    observed sequences of actions, not as proven intent.
  \item \textbf{No general superiority claim.} Any observed advantage for
    event-driven or periodic collection in this study is scoped to the
    tested cap, rate, and interval settings and must not be read as a
    general claim that one architecture dominates the other across all
    monitoring settings or domains.
\end{itemize}

\subsection*{Dual-Use Considerations}
\begin{itemize}[leftmargin=*]
  \item \textbf{Defensive framing.} This work is intended to inform
    \emph{defensive} monitoring/observability design for AI incident
    response (i.e.\ helping defenders retain the evidence they need).
  \item \textbf{Potential misuse vector.} The same evidence-retention
    analysis could in principle be used by an adversarial actor to identify
    which monitoring configurations retain the \emph{least} evidence of a
    given behavior, i.e.\ to time or structure actions to evade a specific
    monitoring cadence. \todo{state directly whether/why the team judges the
    defensive value to outweigh this risk, and whether any specific
    technique in the repository (e.g.\ exact evasion-timing parameters) was
    withheld from public release for this reason.}
  \item \textbf{No exploit code released.} \todo{confirm and state whether
    the public repository (if published) omits anything that would function
    as a ready-made evasion-timing recipe against a real monitoring system.}
\end{itemize}

% ==========================================================================
\section*{References}
% ==========================================================================
% \todo Replace with real citations (related work, any external methodology
% referenced). A plain bibliography is used here instead of BibTeX/biblatex
% to keep this template dependency-free; swap in \bibliography{refs} if the
% team prefers a .bib file.

\begin{thebibliography}{9}
\bibitem{todo1} \todo{Author, ``Title,'' Venue, Year.}
\bibitem{todo2} \todo{Author, ``Title,'' Venue, Year.}
\end{thebibliography}

\end{document}
